% GENERATED FILE. Edit canonical lessons, then run npm run generate:books.
% canonical-source-hash: fnv1a64:cefb28a002e02ba7
% canonical-lessons: MR-C340-rahasya, MR-C340-tatparya, MR-C340-vilaksan, MR-C340-thararak, MR-C340-jhagadne

\chapter{Mysteries}
\label{ch:mr-mysteries}
\hlchaptermodality{340}

\begin{chapteropening}
\textbf{By the end of this chapter:} \emph{I can tell a thrilling story with a mystery and a moral.}

Complete the last lesson of chapter 340: I can tell a thrilling story with a mystery and a moral.
\end{chapteropening}

\section[\texorpdfstring{\mr{रहस्य} (rahasya)}{रहस्य (rahasya)}]{\mr{रहस्य} (rahasya) — a mystery; a secret}

\label{lesson:MR-C340-rahasya}



\begin{quote}
Before the new one: say the Marathi for a holy man, then the Marathi for a moneylender.
\end{quote}

\subsection*{You'll want to know: \mr{रहस्य}}
\textbf{\mr{रहस्य}} — \emph{rahasya} — \textquotedblleft{}a mystery; a secret\textquotedblright{}.

\begin{quote}
\mr{त्या} \mr{जुन्या} \mr{वाड्याचं} \mr{रहस्य} \mr{कोणालाच} \mr{माहीत} \mr{नाही}.
\end{quote}

\begin{quote}
\emph{tyā junyā vāḍyācaṁ rahasya koṇālāc māhīt nāhī.}
\end{quote}

\begin{quote}
Nobody knows the secret of that old mansion.
\end{quote}

Neuter.

\mr{रहस्यकथा}, a mystery story.

\begin{cousinweb}
\textbf{\mr{रहस्य}} (\emph{rahasya}) is a tatsama, taken from Sanskrit \emph{rahasya}, 'secret, private; a mystery' (Monier-Williams).
\end{cousinweb}

\begin{grammarlens}[title={What you've built}]
One more word for telling what happened.
\end{grammarlens}

\subsection*{Your turn}
\begin{itemize}
\raggedright
  \item \emph{Say it:} \emph{rahasya}
  \item \emph{Say it:} \emph{rahasya}, once more
  \item \emph{Say it:} \emph{tyā junyā vāḍyācaṁ rahasya koṇālāc māhīt nāhī.} — Nobody knows the secret of that old mansion.
  \item \emph{Recall:} say the Marathi for a holy man, then the Marathi for a moneylender, then say \emph{rahasya} again
\end{itemize}

\begin{tcolorbox}[breakable,colback=teal!4,colframe=teal!35!black,title={Before you move on}]
What does \mr{रहस्य} mean? (\textquotedblleft{}A mystery; a secret\textquotedblright{}.) Say it once more.
\end{tcolorbox}
\section[\texorpdfstring{\mr{तात्पर्य} (tātparya)}{तात्पर्य (tātparya)}]{\mr{तात्पर्य} (tātparya) — the moral; the gist}

\label{lesson:MR-C340-tatparya}



\begin{quote}
Before the new one: say the Marathi for a moneylender, then the Marathi for a mystery.
\end{quote}

\subsection*{You'll want to know: \mr{तात्पर्य}}
\textbf{\mr{तात्पर्य}} — \emph{tātparya} — \textquotedblleft{}the moral; the gist\textquotedblright{}.

\begin{quote}
\mr{या} \mr{गोष्टीचं} \mr{तात्पर्य} \mr{असं} \mr{आहे} \mr{की} \mr{खोटं} \mr{बोलू} \mr{नये}.
\end{quote}

\begin{quote}
\emph{yā goṣṭīcaṁ tātparya asaṁ āhe kī khoṭaṁ bolū naye.}
\end{quote}

\begin{quote}
The moral of this story is that one should not lie.
\end{quote}

Neuter.

Children's stories end with \mr{तात्पर्य}, the moral.

\begin{cousinweb}
\textbf{\mr{तात्पर्य}} (\emph{tātparya}) is a tatsama, taken from Sanskrit \emph{tātparya}, 'aim, purport' (Monier-Williams).
\end{cousinweb}

\begin{grammarlens}[title={What you've built}]
One more word for telling what happened.
\end{grammarlens}

\subsection*{Your turn}
\begin{itemize}
\raggedright
  \item \emph{Say it:} \emph{tātparya}
  \item \emph{Say it:} \emph{tātparya}, once more
  \item \emph{Say it:} \emph{yā goṣṭīcaṁ tātparya asaṁ āhe kī khoṭaṁ bolū naye.} — The moral of this story is that one should not lie.
  \item \emph{Recall:} say the Marathi for a moneylender, then the Marathi for a mystery, then say \emph{tātparya} again
\end{itemize}

\begin{tcolorbox}[breakable,colback=teal!4,colframe=teal!35!black,title={Before you move on}]
What does \mr{तात्पर्य} mean? (\textquotedblleft{}The moral; the gist\textquotedblright{}.) Say it once more.
\end{tcolorbox}
\section[\texorpdfstring{\mr{विलक्षण} (vilakṣaṇ)}{विलक्षण (vilakṣaṇ)}]{\mr{विलक्षण} (vilakṣaṇ) — extraordinary; strange}

\label{lesson:MR-C340-vilaksan}



\begin{quote}
Before the new one: say the Marathi for a mystery, then the Marathi for the moral.
\end{quote}

\subsection*{You'll want to know: \mr{विलक्षण}}
\textbf{\mr{विलक्षण}} — \emph{vilakṣaṇ} — \textquotedblleft{}extraordinary; strange\textquotedblright{}.

\begin{quote}
\mr{त्या} \mr{रात्री} \mr{एक} \mr{विलक्षण} \mr{गोष्ट} \mr{घडली}.
\end{quote}

\begin{quote}
\emph{tyā rātrī ek vilakṣaṇ goṣṭa ghaḍlī.}
\end{quote}

\begin{quote}
Something extraordinary happened that night.
\end{quote}

\mr{विलक्षण} keeps one form.

\begin{cousinweb}
\textbf{\mr{विलक्षण}} (\emph{vilakṣaṇ}) is a tatsama, taken from Sanskrit \emph{vilakṣaṇa}, 'having different marks, unlike others' (Monier-Williams).
\end{cousinweb}

\begin{grammarlens}[title={What you've built}]
One more word for telling what happened.
\end{grammarlens}

\subsection*{Your turn}
\begin{itemize}
\raggedright
  \item \emph{Say it:} \emph{vilakṣaṇ}
  \item \emph{Say it:} \emph{vilakṣaṇ}, once more
  \item \emph{Say it:} \emph{tyā rātrī ek vilakṣaṇ goṣṭa ghaḍlī.} — Something extraordinary happened that night.
  \item \emph{Recall:} say the Marathi for a mystery, then the Marathi for the moral, then say \emph{vilakṣaṇ} again
\end{itemize}

\begin{tcolorbox}[breakable,colback=teal!4,colframe=teal!35!black,title={Before you move on}]
What does \mr{विलक्षण} mean? (\textquotedblleft{}Extraordinary; strange\textquotedblright{}.) Say it once more.
\end{tcolorbox}
\section[\texorpdfstring{\mr{थरारक} (tharārak)}{थरारक (tharārak)}]{\mr{थरारक} (tharārak) — thrilling}

\label{lesson:MR-C340-thararak}



\begin{quote}
Before the new one: say the Marathi for the moral, then the Marathi for extraordinary.
\end{quote}

\subsection*{You'll want to know: \mr{थरारक}}
\textbf{\mr{थरारक}} — \emph{tharārak} — \textquotedblleft{}thrilling\textquotedblright{}.

\begin{quote}
\mr{त्या} \mr{चित्रपटाचा} \mr{शेवट} \mr{फार} \mr{थरारक} \mr{होता}.
\end{quote}

\begin{quote}
\emph{tyā citrapaṭācā śevaṭ phār tharārak hotā.}
\end{quote}

\begin{quote}
The ending of that film was very thrilling.
\end{quote}

\mr{थरारक} keeps one form; \mr{थरार}, a thrill.

\begin{cousinweb}
\textbf{\mr{थरारक}} (\emph{tharārak}) is built on \emph{tharār}, 'a thrill, a shiver'; CDIAL lists Marathi \emph{tharārṇe}, 'to tremble', under \emph{*thar}, 'tremble' (CDIAL 6092).
\end{cousinweb}

\begin{grammarlens}[title={What you've built}]
One more word for telling what happened.
\end{grammarlens}

\subsection*{Your turn}
\begin{itemize}
\raggedright
  \item \emph{Say it:} \emph{tharārak}
  \item \emph{Say it:} \emph{tharārak}, once more
  \item \emph{Say it:} \emph{tyā citrapaṭācā śevaṭ phār tharārak hotā.} — The ending of that film was very thrilling.
  \item \emph{Recall:} say the Marathi for the moral, then the Marathi for extraordinary, then say \emph{tharārak} again
\end{itemize}

\begin{tcolorbox}[breakable,colback=teal!4,colframe=teal!35!black,title={Before you move on}]
What does \mr{थरारक} mean? (\textquotedblleft{}Thrilling\textquotedblright{}.) Say it once more.
\end{tcolorbox}
\section[\texorpdfstring{\mr{झगडणे} (jhagaḍṇe)}{झगडणे (jhagaḍṇe)}]{\mr{झगडणे} (jhagaḍṇe) — to wrestle with; to struggle}

\label{lesson:MR-C340-jhagadne}



\begin{quote}
Before the new one: say the Marathi for extraordinary, then the Marathi for thrilling.
\end{quote}

\subsection*{You'll want to know: \mr{झगडणे}}
\textbf{\mr{झगडणे}} — \emph{jhagaḍṇe} — \textquotedblleft{}to wrestle with; to struggle\textquotedblright{}.

\begin{quote}
\mr{नायक} \mr{शेवटपर्यंत} \mr{संकटांशी} \mr{झगडला}.
\end{quote}

\begin{quote}
\emph{nāyak śevaṭparyant saṅkaṭāṁśī jhagaḍlā.}
\end{quote}

\begin{quote}
The hero wrestled with troubles to the very end.
\end{quote}

Past (he …): \mr{तो} \mr{झगडला}, \emph{to jhagaḍlā}.

The opponent takes -\mr{शी}: \mr{संकटांशी} \mr{झगडणे}; also to quarrel.

\begin{cousinweb}
\textbf{\mr{झगडणे}} (\emph{jhagaḍṇe}) is inherited: CDIAL lists Marathi \emph{jhagaḍṇe}, 'to quarrel', under \emph{*jhagaḍ}, 'quarrel' (CDIAL 5321).
\end{cousinweb}

\begin{grammarlens}[title={What you've built}]
One more word for telling what happened.
\end{grammarlens}

\subsection*{Your turn}
\begin{itemize}
\raggedright
  \item \emph{Say it:} \emph{jhagaḍṇe}
  \item \emph{Say it:} \emph{jhagaḍṇe}, once more
  \item \emph{Say it:} \emph{nāyak śevaṭparyant saṅkaṭāṁśī jhagaḍlā.} — The hero wrestled with troubles to the very end.
  \item \emph{Recall:} say the Marathi for extraordinary, then the Marathi for thrilling, then say \emph{jhagaḍṇe} again
\end{itemize}

\begin{tcolorbox}[breakable,colback=teal!4,colframe=teal!35!black,title={Before you move on}]
What does \mr{झगडणे} mean? (\textquotedblleft{}To wrestle with; to struggle\textquotedblright{}.) Say it once more.
\end{tcolorbox}
